\documentclass[10pt,a4paper]{amsart}
\usepackage[margin=19mm]{geometry}
\usepackage{amsmath,amssymb,mathtools}
\usepackage[scr=rsfs,cal=euler]{mathalfa}
\usepackage{xfrac}
\usepackage[unicode,hidelinks,bookmarks=false]{hyperref}
\newcommand{\ie}{i.e.\ }
\newcommand{\cf}{cf.\ }
\newcommand{\resp}{respectively\ }
\newcommand{\Subsection}{\subsection}
\providecommand{\nobreakdash}{\nobreak\hbox{-}}
\setlength{\emergencystretch}{2em}
\multlinegap0pt
\usepackage{bm}
\usepackage{bbm}
\usepackage{dsfont}
\usepackage{xspace}
\usepackage{cancel}
\usepackage{mathtools}
\usepackage{amsthm}
\makeatletter
\@ifundefined{thm}{\newtheorem{thm}{Theorem}}{}
\makeatother
\def \ie{{\it i.e.}}
\def\IC{\mathbb{C}}
\def\IZ{\mathbb{Z}}
\def\df{\mathfrak{d}}
\def\gf{\mathfrak{g}}
\title[Twisted chiral algebras of class $\mathcal{S}$ and mixed Fei]{Twisted chiral algebras of class $\mathcal{S}$ and mixed Feigin-Frenkel gluing}
\author{Christopher Beem, Sujay Nair}
\date{Source: arXiv:2201.13435v1 (CC BY 4.0)}
\begin{document}
\maketitle

\noindent\emph{Source and licence.} Twisted chiral algebras of class $\mathcal{S}$ and mixed Feigin-Frenkel gluing. Version arXiv:2201.13435v1 (CC BY 4.0), distributed under the Creative Commons Attribution 4.0 International licence, as recorded in the arXiv licence metadata for this exact version and shown on the abstract page https://arxiv.org/abs/2201.13435v1. Full rights notice: licenses/arxiv-2201.13435-CC-BY-4.0.txt.\par\smallskip

\noindent\textit{Editorial scope.} Counted unit: the Thm whose source label is \texttt{None} in the article (source0.tex lines 2777--2779, source label \texttt{None}), delivered together with the article's own complete original proof of it (source0.tex lines 2781--2866, one \texttt{proof} environment of 86 lines). No mathematics is written, repaired or re-derived: statement and proof are the article's own text line for line. Required context retained uncounted: none. Every definition, display and label of those lines is retained unchanged. Omitted from this package and not redistributed: 1--2776, 2780--2780, 2867--3373. Nothing inside a retained excerpt is omitted or altered: every retained body is delivered exactly as the article's lines are written, so no omission receipt of any kind accompanies this package. Citations: the retained excerpts carry no citation command at all, so no bibliography entry of the article is transcribed and no reference is redistributed. Reference adaptation: none was needed and none was made; every reference inside the delivered unit resolves inside it, so no unresolved reference or citation is produced. Printed slips kept literal: no printed word, symbol, hypothesis or number of the counted statement or of its proof was altered. Layout and class: the article's own class is \texttt{article} with packages including pdf14, fix-cm, ./aux/jheppub, graphicx, amsfonts, amsmath, amssymb, amsthm, tikz, tikz-cd, ./aux/tikzit, verbatim, parskip, array; the wrapper is the lane's pinned \texttt{amsart} editorial wrapper at 10pt on A4, which restates the theorem-like environments the retained text needs and adds the article's own macro definitions that the retained text needs (\texttt{\textbackslash IC}, \texttt{\textbackslash IZ}, \texttt{\textbackslash df}, \texttt{\textbackslash gf}, \texttt{\textbackslash ie}), with extra wrapper packages (bm, bbm, dsfont, xspace, cancel, mathtools). The delivered numbering is the unit's own. Assets: no retained span contains a figure environment, a graphics-inclusion command, a PDF-page inclusion, a table or a TikZ picture; no image file is copied into this package, the package declares no asset dependency and no image file is redistributed with it. Licence: version arXiv:2201.13435v1 of this article is distributed under the Creative Commons Attribution 4.0 International licence, as recorded in the arXiv licence metadata for this exact version and shown on the abstract page https://arxiv.org/abs/2201.13435v1. A full rights notice is delivered at licenses/arxiv-2201.13435-CC-BY-4.0.txt. Characters: every retained excerpt is pure ASCII, so the delivered engine, XeLaTeX with its default Latin Modern text font, reports no missing character.
\begin{thm}
The map $\pi:\mathrm{Fun\, Op}_{G_u}^\lambda \rightarrow \mathrm{Fun\, Op}_{^LG_t}^\lambda$, defined as above, is surjective.
\end{thm}

\begin{proof}
%%%
The action of $\mathfrak{n}_u^\sigma$ on the space of coinvariants can be realised through the symmetrised screening charges
%%%%%%
\begin{equation}
	\tilde{V}_i[\lambda_i +1] = \frac{1}{|\langle \sigma \rangle|} \sum_{\sigma'\in\langle \sigma \rangle} V_{\sigma(i)}[\lambda_{\sigma(i)}+1]~.
\end{equation}
%%%%%%
A polynomial $P\in\mathrm{Fun\, Conn}(\Omega^{-\rho})$ is in the space of invariants, $(\mathrm{Fun\, Conn}(\Omega^{-\rho}))^{\mathfrak{n}_u}$, if and only if it is in the intersection of the kernels of the screening charges,~\ie
%%%%%%
\begin{equation}
	V_i[\lambda_i +1]P \overset{!}= 0\,,\qquad \text{for } i=1,\dots,\mathrm{rk}\,\gf_u~.
\end{equation}
%%%%%%
Similarly, a polynomial $P\in(\mathrm{Fun\, Conn}(\Omega^{-\rho}))_\sigma$ in the space of coinvariants is an $\mathfrak{n}_u^\sigma$ invariant, if and only if it lies in the intersection of the kernels of the symmetrised screening charges $\tilde{V}_i[\lambda_i+1]$.

It is, therefore, sufficient to show that if $P$ is a $\sigma$-coinvariant and lies in the interection of the kernels of the symmetrised screening charges, then it must lie in the intersection of the kernels of the $\mathfrak{n}_u$ screening charges.
We make the following observations. Suppose $P$ lies in the space of coinvariants, then we can realise it as $P\in\IC[\tilde{u}_{i,n}|i=1,\dots,\mathrm{rk}\,\gf_t;\,n<0]$. Therefore,
%%%%%%
\begin{equation}
	\frac{\partial P}{\partial u_{j,n}} = \sigma \left( \frac{\partial P}{\partial u_{j,n}}\right) = \frac{\partial P}{\partial u_{\sigma(j),n}} ~.
\end{equation}
%%%%%%
Furthermore, the $x_{i,m}$ must transform as,
%%%%%%
\begin{equation}
	\sigma(x_{i,m}) = x_{\sigma(i),m}~,
\end{equation}
%%%%%%
which follows straightforwardly from the definition. Combining, we have that
%%%%%%
\begin{equation}\label{eq:sigma_screening_charge}
	\sigma\left(V_i[\lambda_i + 1] P\right) = V_{\sigma(i)}[\lambda_{\sigma(i)}+1]P~.
\end{equation}
%%%%%%
Let $P\in(\mathrm{Fun\, Conn}(\Omega^{-\rho}))_\sigma$ and suppose it lies in the intersection of kernels of the symmetrised operators. Additionally, let us restrict to the case where $\sigma^2=1$ and we only consider the $\IZ_2$ twist---we shall address the $\IZ_3$ twist for $\df_4$ at the end. Now, we have that
%%%%%%
\begin{equation}
	(V_{i}[\lambda_i+1] + V_{\sigma(i)}[\lambda_{\sigma(i)}+1])P = 0~,
\end{equation}
%%%%%%
but from (\ref{eq:sigma_screening_charge}), we must have that
%%%%%%
\begin{equation}
	\sigma\left(V_{i}[\lambda_i+1]P\right)=-V_{i}[\lambda_i+1]P~.
\end{equation}
%%%%%%
We shall now show that for any $P$ in the space of coinvariants, the polynomial $V_i[\lambda_i+1]P$, cannot have eigenvalue $-1$ under $\sigma$.

The image of $P$ under the $i$th screening charge is
%%%%%%
\begin{equation}
	V_i[\lambda_i +1]P = \sum_{n\ge\lambda_i} x_{i,n-\lambda_i}\cdot \left( a_{ij} \frac{\partial P}{\partial u_{j,-n-1}} \right)~,
\end{equation}
%%%%%%
and we have
%%%%%%
\begin{equation}
	\sigma \left( V_i[\lambda_i +1]P \right) = V_{\sigma(i)}[\lambda_{i}+1]P= \sum_{n\ge\lambda_{i}} x_{\sigma(i),n-\lambda_{i}}\cdot \left( a_{ij} \frac{\partial P}{\partial u_{j,-n-1}} \right)~,
\end{equation}
%%%%%%
where we have made use of the fact that the derivatives of $P$ and the highest weight, $\lambda$, are invariant under $\sigma$. From inspection, it is clear that
%%%%%%
\begin{equation}
	V_{\sigma(i)}[\lambda_{\sigma(i)}+1]P\neq- V_i[\lambda_i +1]P ~,
\end{equation}
%%%%%%
unless both are identically zero, as desired.

Now we address the $\IZ_3$ case of $\gf_u=\df_4$ and $\gf_t=\gf_2$. A $\mathfrak{n}_u^\sigma$ invariant $P$ must satisfy the following:
%%%%%%
\begin{equation}
	\begin{split}
	(V_{1}[\lambda_1+1] + V_{3}[\lambda_1+1]+V_4[\lambda_1+1])P=0~,\\
	V_{2}[\lambda_2+1]P=0~.
	\end{split}
\end{equation}
%%%%%%
We have, once more, made use of the fact that the derivatives of $P$ and the highest weight $\lambda$ are invariant under the action of $\IZ_3$. Expanding the first requirement, we have that
%%%%%%
\begin{equation}
	\sum_{n\ge\lambda_1} \left(x_{1,n-\lambda_1}+ x_{3,n-\lambda_1}+x_{4,n-\lambda_i}\right) \left( -\frac{\partial P}{\partial u_{2,-n-1}} + 2 \frac{\partial P}{\partial u_{i,-n-1}} \right) =0~.
\end{equation}
%%%%%%
Once again, this can only hold if each screening charge individually acts as zero.
\end{proof}

\end{document}
